\chapter{Reproducibility record}\label{app:repro}

\section{Code and environment}

The project repository at \url{https://github.com/whynotramaa/fcsg-capstone}
holds all code, configuration files, logs and evaluation outputs. Model code
is in \texttt{models/} and \texttt{src/fcsg\_net/}, training in
\texttt{training/train.py}, and evaluation and plotting in
\texttt{evaluation/}. A configuration file in \texttt{configs/} selects each
model. The four composite runs use \texttt{fcsg.toml},
\texttt{dncnn\_composite.toml}, \texttt{ffdnet\_blind.toml} and
\texttt{ffdnet\_composite.toml}. They differ only in the model and the noise
map.

We ran training, verification and diagnostics on Kaggle with an NVIDIA Tesla
T4, Python 3.12.13, PyTorch 2.10.0 with CUDA 12.8, NumPy 2.0.2 and Pillow
11.3.0. Local machines were only for editing and packaging code.

\section{Where each result comes from}

\begin{table}[htbp]
	\centering
	\caption{Source files for the results in this report.}
	\label{tab:sources}
	\small
	\begin{tabular}{@{}ll@{}}
		\toprule
		Result & File \\
		\midrule
		Verification checks (Table~\ref{tab:gates}) & \texttt{results/milestone\_run\_20260925/} \\
		Gaussian DnCNN (Table~\ref{tab:gaussian})   & \texttt{results/benchmark.csv} \\
		Composite results (Table~\ref{tab:main})     & \texttt{results/phase3/*/evaluation.json} \\
		Per-image results (Appendix~\ref{app:per-image}) & \texttt{results/phase3/per\_image\_scores.csv} \\
		Training curves                              & \texttt{results/phase3/*/train\_log.csv} \\
		Latency, tiling, routing                     & \texttt{results/diagnostics/} \\
		Figures                                      & \texttt{output/pdf/fcsg-report-assets/} \\
		\bottomrule
	\end{tabular}
\end{table}

\section{Artefact identifiers}

Each Kaggle run packed its source code into a bundle and recorded the
bundle's SHA-256 hash. Even after later edits, the hash pins down the exact
code behind each result. Table~\ref{tab:hashes} lists the first sixteen hex
digits of each hash. The source manifests hold the full values.

\begin{table}[htbp]
	\centering
	\caption{SHA-256 identifiers (first 16 hex digits) of the evaluated artefacts.}
	\label{tab:hashes}
	\small
	\begin{tabular}{@{}ll@{}}
		\toprule
		Artefact & SHA-256 prefix \\
		\midrule
		Composite training source bundle & \texttt{85d6e16f72ead013} \\
		Diagnostics source bundle        & \texttt{d473bc10b0ad0c1e} \\
		5,000-pair dataset manifest      & \texttt{7ce7f310bdb4d2d5} \\
		\fcsg{} checkpoint               & \texttt{8db0bb0f3c251ddc} \\
		DnCNN checkpoint                 & \texttt{afb0abfa022ede04} \\
		FFDNet blind checkpoint          & \texttt{981560c18cbea134} \\
		FFDNet oracle checkpoint         & \texttt{ef9560e8ed1ebae0} \\
		\bottomrule
	\end{tabular}
\end{table}

\section{How to reproduce}

\begin{enumerate}
	\item Attach the DIV2K dataset to a Kaggle notebook with a GPU. Run
	\texttt{notebooks/checks.ipynb}, then the bundled milestone notebook
	\texttt{notebooks/phase1\_phase2.ipynb}. Both stop at the first failed
	check.
	\item Build and push one training notebook per configuration with
	\texttt{python notebooks/build\_phase3.py fcsg ffdnet\_blind
	ffdnet\_composite dncnn\_composite}. To continue a run that outlasts one
	session, rebuild it with \texttt{--resume}. It picks up from its last
	checkpoint.
	\item At 200,000 updates, each run evaluates the 100 validation images and
	writes \texttt{evaluation.json} and \texttt{status.json}.
	\item To measure latency, tiling and routing, build the diagnostic notebook
	with \texttt{python notebooks/build\_diagnostics.py}. Draw the figures with
	\texttt{python evaluation/plots.py}.
\end{enumerate}
